% Transcription — fonds Alexandre Grothendieck, shelfmark 81, batch 1 (pages 1-20).
% Established from the facsimile archives/batches/81/batch-01.pdf (a mirror of
% https://grothendieck.umontpellier.fr/81.pdf), per the skill
% transcribe-grothendieck. The batch file carries no cover sheet: PDF page N is
% archivists' page N; the pencilled numbers bottom-left (« 1 », « 4 », « 6 »,
% « 7 », « 8 », « 9 », « 10 », « 11 », « 13 » ... « 20 ») read at the PDF page
% of the same number.
%
% Pass: Opus 5.5 (claude-opus-5-5), 2026-09-23 — first pass, unchecked against the pages by a human.
%
% Fourteen of the twenty pages carry a \page marker. Absent: page 1, a sheet
% otherwise blank inscribed in pencil « 0 Cartes sphériques en général » — a
% cover, whose words head the first section; pages 6-7, an attempted proof of
% Fermat's last theorem (a^q + b^q = c^q, Newton's binomial expansion, with a
% boxed example q = 2, 3-4-5) in a rounded school hand that is not his and
% concerns nothing in this folder; page 8, blank apart from a faint pencilled
% name and dates read through the paper; page 12, a table in his hand that is
% not mathematics, whose latest year, 1981, agrees with the archivists' « [à
% partir de 1981] »; page 13, a
% pink cover inscribed « décompositions n-aires (et découpages) d'espaces
% topologiques », whose words head the second section. He does not paginate.
%
% What the batch is. Three unrelated scraps, then the run the folder is about.
%   2, 4   « Carte sphérique (orientée) avec repère »: in the isotopy
%          category, the same as a "special" map f : P^1 -> P^1 (f(0)=0,
%          f(1)=1, f(∞)=∞, étale over P^1 - {0,1,∞}) whose ramification
%          over C* is only over 1 and simple; as a pair P, Q of coprime
%          polynomials; the question whether P'(z)=0 forces P(z)=Q(z), tried
%          on P(T) = aT^2 + bT. Page 4: Aut(C) acts on special maps through
%          Aut(Q-bar), the coefficients being algebraic.
%   3, 5   Cancelled algebra on the versos of 2 and 4: generators e_i, f_i,
%          η of an algebra with a derivation δ (e_i e_1 = 0, f_i e_1 =
%          e_{i+1}, δη = ±e_{d-2}); on 5, upside down to it, derivatives of
%          f = P/Q. Condensed.
%   9-11   Drawings: discs cut by chords, each with the graph of the chords
%          (vertices) and their crossings (edges) beneath, several marked
%          « ok »; page 10 adds P(X), Hom(X,Y) ⊂ P(X × Y) and a sketch of a
%          correspondence. Described, not redrawn.
%   14-20  « Version I »: a space X cut by closed parts E_i (i ∈ I), each
%          separating X into exactly two components Σ_i^+, Σ_i^-, each E_j in
%          one component of X - E_i. The case card I = 2 (Σ_1^+ ⊂ Σ_2^-,
%          the five-piece partition with Σ_12 = Σ_1^- ∩ Σ_2^-, closures); the
%          partition Π_J^0 for J = I - {i} and the induction: n cuts give n+1
%          classes, each of the form X - ∪_{α∈K} (Σ_α^+ ∪ E_α) with the
%          Σ_α^+ pairwise disjoint; a corollary on the uniqueness of that
%          expression, cut off at the foot of page 20. Page 19 states the
%          lemma whose proof pages 18 and 20 give.
%
% The dating: the folder's own. The group « Géométrie et topologie
% combinatoire » (69-89) holds it.
\documentclass[11pt,a4paper]{article}
% Le préambule commun à toutes les transcriptions du fonds Grothendieck.
%
% Il tient en une idée : **l'incertitude se compose, elle ne se résout pas.**
% Une transcription qui lisse un mot douteux a détruit la seule chose pour
% laquelle on la faisait — un lecteur doit pouvoir savoir, à chaque endroit,
% ce qui était sur le feuillet et ce qui a été ajouté. Les six macros qui
% suivent sont donc l'appareil critique, et non de la décoration.
%
% Les mêmes commandes sont reconnues par `scripts/render.mjs`, qui produit la
% vue de lecture du site. Une macro ajoutée ici doit l'être aussi là-bas,
% faute de quoi le rendu échouera bruyamment — ce qui est voulu : un rendu
% silencieusement faux est pire qu'un rendu absent.

\ProvidesPackage{grothendieck}[2026/08/08 Transcriptions du fonds Grothendieck]

\usepackage{amsmath,amssymb,amsthm}
\usepackage{mathrsfs}
\usepackage{stmaryrd}
% Les diagrammes commutatifs sont partout dans ce fonds : c'est la langue
% même de la géométrie algébrique de Grothendieck, et une transcription qui
% les rendrait en prose ne transcrirait rien.
\usepackage{tikz-cd}
% Français par défaut — c'est la langue des feuillets — mais l'anglais est
% chargé aussi : la traduction et le résumé sont des documents anglais, et la
% typographie française (espaces avant les deux-points) y serait une faute.
% Ces documents basculent par \selectlanguage{english} en tête de corps.
\usepackage[english,french]{babel}
\usepackage{fontspec}
\usepackage{geometry}
\geometry{a4paper,margin=25mm,marginparwidth=28mm}
\usepackage{marginnote}
\usepackage{xcolor}
% ulem, not soul. `soul`'s \st and \ul are built on a character-by-character
% scan that XeLaTeX's font machinery breaks: the document fails with an
% undefined control sequence rather than degrading. ulem does the same two jobs
% and is Unicode-engine safe. `normalem` stops it hijacking \emph, which the
% transcriptions use for Grothendieck's own emphasis.
\usepackage[normalem]{ulem}
\usepackage{hyperref}
\hypersetup{colorlinks=true,linkcolor=black,urlcolor=black,pdfborder={0 0 0}}

% --- Métadonnées du lot -----------------------------------------------------
% Reprises telles quelles par le script de rendu, qui les affiche en tête de
% la vue de lecture. Elles fixent ce que le fichier prétend couvrir : sans
% elles, une transcription flotte sans point d'attache dans un fonds de seize
% mille feuillets.

\newcommand{\folder}[1]{\gdef\grothFolder{#1}}
\newcommand{\batch}[1]{\gdef\grothBatch{#1}}
\newcommand{\leaves}[2]{\gdef\grothFirst{#1}\gdef\grothLast{#2}}
\newcommand{\foldertitle}[1]{\gdef\grothFoldertitle{#1}}
\gdef\grothFolder{?}\gdef\grothBatch{?}\gdef\grothFirst{?}\gdef\grothLast{?}\gdef\grothFoldertitle{}

% --- Le résumé --------------------------------------------------------------
%
% Il ouvre la lecture modernisée, sous le titre « Résumé », et il est composé
% en petit. Ce n'est pas un ornement : c'est le seul passage du document qui
% ne soit pas des mathématiques, et un lecteur doit pouvoir le reconnaître
% avant de l'avoir lu — puis le sauter s'il connaît déjà le sujet, ou s'y
% arrêter s'il ne le connaît pas. Le filet qui le suit marque la reprise du
% corps.

\newenvironment{resume}
  {\small\setlength{\parskip}{0.45em}}
  {\par\medskip\hrule\medskip}

% --- Mots-clés ---------------------------------------------------------------
%
% En anglais, en clôture du résumé de la lecture modernisée : le vocabulaire
% moderne sous lequel chercher ce que les feuillets construisent. Le manifeste
% du site (`npm run manifest`) les extrait pour étiqueter la cote entière —
% cette ligne est la source unique des tags, il n'y a pas de fichier à côté.
\newcommand{\keywords}[1]{\par\smallskip\noindent
  {\footnotesize\textsc{Keywords} --- \textit{#1}}\par}

% --- L'appareil critique ----------------------------------------------------

% Le numéro de feuillet, en marge. C'est l'ancre qui permet de régler un
% désaccord sur une formule en regardant *un* feuillet plutôt qu'une chemise
% de six cents. Le site s'en sert aussi pour tourner les pages du fac-similé
% quand on fait défiler la transcription.
\newcommand{\leaf}[1]{%
  \marginnote{\footnotesize\color{gray!70}\textsf{#1}}%
  \label{leaf:#1}\ignorespaces}

% Illisible : on ne devine pas. Un mot inventé qui se lit comme les autres est
% le pire résultat possible.
\newcommand{\ill}{\textcolor{red!60!black}{[\,\ldots\,]}}

% Lecture douteuse : le mot est donné, et signalé comme incertain.
\newcommand{\uncertain}[1]{\uline{#1}}

% Ajout de l'éditeur — une lettre manquante, un mot sous-entendu.
\newcommand{\add}[1]{\textcolor{blue!50!black}{[#1]}}

% Biffé par Grothendieck lui-même. Ce qu'il a rayé fait partie du document :
% une rature dit souvent où la pensée a changé de direction.
\newcommand{\struck}[1]{\sout{#1}}

% Note du transcripteur, sur l'état matériel du feuillet.
\newcommand{\note}[1]{\par\smallskip{\small\color{gray!60!black}\textit{[#1]}}\par\smallskip}

% Note marginale de Grothendieck, distincte du corps du texte.
\newcommand{\marginal}[1]{\par\smallskip\noindent\hspace{1em}%
  {\small\color{gray!60!black}#1}\par\smallskip}

% --- Titre ------------------------------------------------------------------

\AtBeginDocument{%
  \begin{center}
    {\large\bfseries Fonds Alexandre Grothendieck --- cote n\textsuperscript{o}~\grothFolder}\\[2pt]
    {\itshape\grothFoldertitle}\\[4pt]
    {\small feuillets \grothFirst--\grothLast\ (lot \grothBatch)}\\[2pt]
    {\footnotesize\color{gray!60!black}Transcription établie d'après le fac-similé numérisé
     par l'Université de Montpellier.\\
     Les lectures incertaines sont soulignées, les illisibles notées [\,\ldots\,],
     les ajouts entre crochets.}
  \end{center}
  \vspace{1em}}

\endinput


\folder{81}
\batch{1}
\pages{1}{20}
\dating{[à partir de 1981]}
\watermark{Édition de démonstration}
\foldertitle{Immersions du disque et de la circonférence : notes manuscrites (s.d.).}

\begin{document}

\section{Cartes sphériques en général}
\note{ces mots sont de sa main, au crayon, en haut à droite du feuillet 1, par
ailleurs vierge, qui ne reçoit pas de numéro de page ; ils y sont précédés
d'un « 0 »}

\page{2}
Carte sphérique \add{(orientée)} avec repère (dans cat.\ isotopiques)

$\simeq$ morphismes $\mathbb{P}^1_{\mathbb{C}} \xrightarrow{\ f\ }
\mathbb{P}^1_{\mathbb{C}}$ “spécial” i.e.
$f(0)=0$, $f(1)=1$, $f(\infty)=\infty$. $f$ étale en dessus de
$\mathbb{P}^1 - \{0,1,\infty\}$

$\simeq$ fonctions rationnelles $f(z)$ telles que $f(0)=0$, $f(1)=1$,
$f(\infty)=\infty$ et telles que pour $z\in\mathbb{C}$ \struck{telle que}
\uncertain{avec} $f(z)\notin\{0,\infty\}$, $f'(z)=0$, on ait $f(z)=1$
\struck{\ill{}} \add{et de plus} $f''(z)\neq 0$ (i.e.\ la ramification de $f$
en un $z$ est $2$)….

$\simeq$ [$f=P/Q$ !] couples $P,Q$ (de polynômes $\in\mathbb{C}[z]$) premiers
entre eux \add{(i.e.\ sans zéros communs)} tels que $P(0)=0$ (donc
$Q(0)\neq 0$) $P(1)=Q(1)$ (donc $\neq 0$), $\deg Q<\deg P$ (ces conditions
expriment resp.\ ! $f(0)=0$ $f(1)=1$, $f(\infty)=\infty$) normalisés (par
exemple) par $P(1)=Q(1)=1$, et tels que pour $z\in\mathbb{C}$ avec
[$f(z)\notin\{0,\infty\}$, $f'(z)=0$, i.e.] \add{$P(z)\neq 0$,} $Q(z)\neq 0$,
\[
\frac{P'(z)}{P(z)} - \frac{Q'(z)}{Q(z)} = 0 ,
\]
on ait $P(z)=Q(z)$, et de plus
$\dfrac{P''(z)}{P(z)} - \dfrac{Q''(z)}{Q(z)} \neq 0$ i.e.\ $P''(z)\neq Q''(z)$…
\marginal{on pourrait \uncertain{aussi} normaliser \ill{} $Q(0)=1$ \ill{}
\ill{}}

[NB on a
\[
Q(z)P'(z) - P(z)Q'(z) = (Q(z)-P(z))\,P'(z) + P(z)\,(P'(z)-Q'(z))
\]
\note{la ligne porte d'abord, biffé, le second membre
\struck{$Q(z)\,[P'(z)-P(z)] + P(z)\,(Q(z)-Q'(z))$} ; le bon est écrit au-dessus}
pour $P(z), Q(z)\neq 0$, donc si cette quantité est nulle (i.e.\ $f'(z)=0$), alors $P(z)=Q(z)$ \struck{\ill{}} i.e.\ $Q(z)-P(z)=0$ implique
$P'(z)=Q'(z)$ et inversement si $P'(z)=Q'(z)$ on a $(Q(z)-P(z))\,P'(z)=0$
donc si $P'(z)\neq 0$, on a $P(z)=Q(z)$ \add{donc $P'(z)=Q'(z)$}. Mais
peut-on avoir $P'(z)=0$ (d'où aussi $Q'(z)=0$) \ill{} sans avoir
$P(z)=Q(z)$ ?

\struck{Prenons $Q(T)=$} Prenons p.ex.\ $Q=1$, les conditions sur $P$
deviennent \add{(conséquences des deux précédentes)} $P(0)=0$, $P(1)=1$,
$\deg P\geq 1$, donc \struck{$P(T) =$} $P(T)=T\,P_0(T)$ où on a
\add{(pour $P(z)\neq 0$)} $P_0(1)=1$, la question devient si on peut avoir
$P'(z)=0$ sans avoir $P(z)=1$ \uncertain{contre-ex.}\ p.ex.\ si
$P(T)=T(aT+b)=aT^2+bT$, $a+b=1$, on a
\marginal{$P'(T)=2aT+b$, $z=-\frac{b}{2a}$,
$P\bigl(-\frac{b}{2a}\bigr) = -\frac{b}{2a}\bigl(-\frac{b}{2}+b\bigr) =
-\frac{b^2}{4a}$ ; prenons une valeur \ill{} ($a\neq 0$, $b\neq 0$) \ill{}}
\note{la phrase s'arrête au bord inférieur du feuillet, déchiré ; le calcul
se poursuit, au crayon, le long de la marge gauche, écrit de bas en haut, et
n'y est lu qu'en partie}

\page{3}
\note{feuillet barré de deux grands traits en croix, et sans rapport visible
avec les cartes de la page précédente : relations dans une algèbre engendrée
par des $e_i$, $f_i$ et $\eta$, avec une dérivation $\delta$. Transcription
condensée, des formules lisibles seulement}

$f_{d-2}\,e_1 = -e_{d-2}\,f_1$ ($=\eta_{d-1} = \eta$)

$\delta(\eta) = e_{d-2}$

Encadré, pour $1\leq i\leq d-2$ :
\[
e_i e_1 = 0, \qquad
f_i f_1 = \begin{cases} 0 & \text{si } i \text{ impair}\\
f_{i+1} & \text{si } i \text{ pair}\end{cases}, \qquad
f_i e_1 = e_{i+1}, \qquad
e_i f_1 = \begin{cases} -e_{i+1} & \text{si } i \text{ impair}\\
0 & \text{si } i \text{ pair}\end{cases}
\]

$\eta = e_{d-1}$, $\delta(e_{d-1}) = e_{d-2}$ ; $x f_1 = f_1 x = 0$ si
$\deg x$ pair ;
\[
(f_1+\lambda e_1)(f_1+\lambda e_1) = f_1f_1 + \lambda^2 e_1e_1 +
\lambda(f_1e_1 + e_1f_1), \qquad e_2 = f_1e_1, \quad e_3 = f_2e_1 .
\]

À droite, pour $1\leq i\leq d-2$ :
\[
e_ie_1 = 0,\qquad f_if_1 = \begin{cases} 0 & i \text{ impair}\\
f_{i+1} & i \text{ pair}\end{cases},\qquad
f_ie_1 = \eta_{i+1} = \alpha e_{i+1} + \beta f_{i+1},
\]
\[
e_if_1 = \begin{cases} -\eta_{i+1} = -\alpha e_{i+1} - \beta f_{i+1} &
i \text{ impair}\\
-\eta_{i+1} + e_{i+1} = (1-\alpha)e_{i+1} - \beta f_{i+1} & i \text{ pair}
\end{cases}
\]
$f_1e_1 = -e_1f_1 = \eta_2 = \alpha e_2 + \beta f_2$ ; pour
$1\leq i\leq d-3$,
\[
(e_if_1)e_1 = \begin{cases} -\eta_{i+1}e_1 = -\beta f_{i+1}e_1 =
-\beta\eta_{i+2} & i \text{ impair}\\
-\eta_{i+1}e_1 + \ill{} = -\beta\eta_{i+2} & i \text{ pair}\end{cases}
\]
\[
(e_if_1)e_1 = e_i(f_1e_1) = -e_i(e_1f_1) = -(e_ie_1)f_1 = 0
\]
\note{sur la page, la seconde chaîne d'égalités descend verticalement sous
le premier membre}
donc $\beta=0$, $\alpha=\pm 1$, donc $\delta\eta = \pm e_{d-2}$ ;
$\eta = f_{d-2}e_1$, $\delta\eta = \delta(f_{d-2})e_1 = f_{d-3}e_1$.

\page{4}
\emph{NB} Si $\sigma$ est un automorphisme du corps des complexes, et si
$f$ est “spécial”, alors $f^{\sigma}$ est également spécial — donc
$\mathrm{Aut}(\mathbb{C})$ opère sur l'ens.\ des $f$ spéciaux. Mais il opère
par l'intermédiaire de $\mathrm{Aut}(\overline{\mathbb{Q}})$, car si
$f=P/Q$ spécial, les coeff.\ de $P$, $Q$ sont des \emph{nbs algébriques}. Il
revient au même de dire que les zéros de $P$, $Q$ sont des nbs algébriques,
i.e.\ le diviseur de $f$ est
\note{la phrase s'arrête là, au milieu du feuillet}

\page{5}
\note{feuillet de brouillon à deux couches, écrites tête-bêche, l'une barrée
d'une croix ; transcription condensée des formules lisibles}

Couche barrée, suite des calculs de la page 3 :
\[
\delta(e_1e_{d-2}) = -e_1\,\delta(e_{d-2}) = -e_1e_{d-3}, \qquad
\delta(a\eta) = a\,\delta\eta = a\alpha e_{d-2} + a\beta f_{d-2},
\]
\[
\delta(e_1f_{d-2}) = -e_1f_{d-3} = b\alpha e_{d-2} + b\beta f_{d-2},\qquad
\delta(f_1e_{d-2}) = \delta(f_1)e_{d-2} - f_1\delta(e_{d-2}) =
e_{d-2} - f_1e_{d-3},
\]
\[
\delta(c\eta) = c\,\delta\eta = c\alpha e_{d-2} + c\beta f_{d-2} ;
\]
on en tire $e_1e_{d-3}$, $f_1e_{d-3}$, $f_1f_{d-3}$ comme combinaisons
linéaires de $e_{d-1}$ et $f_{d-1}$, à coeff.\ \uncertain{polynômes} en
\ill{} $\alpha$, $\beta$, $a$, $b$, $c$, $d$ ;
$\delta(e_1x_i) = -e_1\,\delta(x_i) = -e_1x_{i-1}$ ($2\leq i$) ;
$\delta(\alpha e_{i+1} + \beta f_{i+1}) = \ill{}$ ;
$\det\begin{pmatrix} a & c\\ b & d\end{pmatrix} = ad-bc$.

Couche tête-bêche, sur $f = P/Q$ :
\[
\Bigl(\frac{f}{g}\Bigr)' = \frac{f'g-g'f}{g^2}, \qquad
(xy)'' = x''y + 2x'y' + xy'', \qquad
\Bigl(\frac{P}{Q}\Bigr)' = \frac{P'}{Q} - P\frac{Q'}{Q^2} =
\frac{P'Q-PQ'}{Q^2},
\]
\[
\Bigl(\frac{P}{Q}\Bigr)'' = \frac{P''}{Q} - 2P'\frac{Q'}{Q^2} + \ill{}, \qquad
\frac{P''Q^2 - Q''PQ - 2P'QQ' + 2PQ'^2}{Q^3} .
\]

\page{9}
\note{feuillet de dessins, sans texte. Trois rangées de disques coupés par
des cordes (droites ou arcs), parfois concourantes en un point marqué, une
région parfois hachurée ; sous chaque disque ou groupe de disques réuni par
une accolade, un petit graphe : un point, deux points, un segment, trois
points, un segment et un point, un chemin à trois sommets, un triangle —
apparemment le graphe dont les sommets sont les cordes et les arêtes leurs
croisements. Plusieurs dessins sont marqués « ok » de sa main (le disque à
trois cordes concourantes au-dessus d'un segment, le disque au triangle
hachuré au-dessus du triangle, avec en variante « ou » un disque à cordes
concourantes) ; sous une rangée de disques à quatre cordes parallèles, un
« ? ». En bas à droite, des lignes brisées épaisses}

\page{10}
$\mathbb{P}(X)$ ; $O\subset\mathbb{P}(X)$
\note{à côté, un dessin : une droite verticale traversant un petit disque
aux rayons marqués, d'où partent des boucles en pétales}
\[
\mathbb{P}(X\times Z) \overset{?}{\simeq}
\mathrm{Hom}(Z\times X,\ \mathcal{E}) =
\mathrm{Hom}(X,\ \underline{\mathrm{Hom}}(Z,\mathcal{E}))
\]
\note{dans le second membre, une première lettre surchargée en $Z$ ; avant
le troisième, \struck{$\underline{\mathrm{Hom}}(X)$}, biffé}
\[
\mathrm{Hom}(\mathcal{E},\ \underline{\mathrm{Hom}}(X,Z)) =
\mathbb{P}(\underline{\mathrm{Hom}}(X,Z))
\]
\note{la première lettre de cette ligne, lue $\mathcal{E}$, est douteuse}
\[
\underline{\mathrm{Hom}}(X,Y)\subset\mathbb{P}(X\times Y), \qquad
\underline{\mathrm{Hom}}(Z\times X,\ Y)\subset\mathbb{P}(Z\times X\times Y)
\]
\note{la lettre notée $\mathcal{E}$ est une grande capitale bouclée, que la
page ne définit pas. À droite, un croquis de flèches : $P\times X\times Y
\supset H$, avec $H\to P\times X$ (flèche oblique) et une flèche verticale
$P\times X\times Y\to P\times X$, puis $P\times X\to P$ ; à côté,
$U\times_V U$ et $V\leftarrow U$ au-dessus d'un $P$. En bas, un polygone et
des lignes brisées épaisses, et une ellipse traversée de courbes}

\page{11}
\note{feuillet de dessins de la même famille que la page 9 : disques coupés
de cordes, une région triangulaire ou quadrangulaire hachurée, trois marqués
« OK », des variantes introduites par « ou » ; sous eux les graphes : un
triangle, un point, un tripode épais en « Y », un carré, un triangle
surmonté d'un segment, un rectangle coupé d'une diagonale, un triangle à
point intérieur. À droite, deux axes fléchés en équerre et, dans l'angle,
des courbes qui se croisent et deux disques coupés d'arcs}

\section{Décompositions \emph{n}-aires (et découpages) d'espaces topologiques}
\note{ces mots sont de sa main, au crayon, en haut à droite du feuillet 13,
une chemise rose par ailleurs vierge qui ne reçoit pas de numéro de page ;
« et découpages » y est ajouté en interligne}

\page{14}
\struck{V} \emph{Version I}

Soient $X$ un espace topologique, $(E_i)_{i\in I}$ un ens.\ de parties
\add{fermées} de $X$ (\uncertain{indexé} \uncertain{indiciellement}), telles
que l'on ait
\begin{itemize}
  \item[(a)] $\forall i\in I$, $\operatorname{card}\pi_0(X\setminus E_i)=2$
  \item[(b)] $\forall (i,j)\in I\times I$, $i\neq j$, $E_j$ est contenu dans
  l'une des deux composantes connexes de $X\setminus E_i$.
\end{itemize}
\marginal{NB les comp.\ conn.\ $\Sigma_i^+$, $\Sigma_i^-$ de
$X\setminus E_i$ sont \uncertain{ouvertes} dans $X$}
NB cette condition (b) est vérifiée automatiquement \ill{} si les $E_i$
($i\in I$) sont connexes, dès lors que les $E_i$ sont supposés disjoints
mutuellement.

On se propose d'étudier le “découpage” de $X$ par la famille $(E_i)$.
\marginal{Supposons de plus que pour \ill{} comp.\ conn.\ $\Sigma_i$ de
$X\setminus E_i$, on ait $\overline{\Sigma}_i\cap E_i\neq\emptyset$ i.e.\ $\overline{\Sigma}_i\neq\Sigma_i$, i.e.\ $\Sigma_i$ \uncertain{n'est pas}
fermé dans $X$ (ce qui est le cas si $X$ est connexe)}

\emph{Cas $\operatorname{card}(I)=2$} Soit $\Sigma_i^+$ la comp.\ conn.\ de $X\setminus E_i$ qui ne contient pas $E_j$ ($j\neq i$), alors on a (si
$I=\{1,2\}$)
\note{à gauche, un dessin : un disque coupé de deux cordes horizontales
$E_1$ (en haut) et $E_2$ (en bas) ; au-dessus de $E_1$, la calotte
$\Sigma_1^+$, et une accolade à droite réunit tout ce qui est sous $E_1$ :
$\Sigma_1^-$ ; sous $E_2$, la calotte $\Sigma_2^+$, et une accolade à gauche
réunit tout ce qui est au-dessus de $E_2$ : $\Sigma_2^-$}
\begin{itemize}
  \item[(1)] $\Sigma_1^+\cap\Sigma_2^+=\emptyset$ ce qui
  \struck{résulte de} \struck{$\Sigma_1^-\supset\Sigma_2^+$} équivaut à :
  \struck{$\Sigma_2^-\supset\Sigma_1^+$} $\Sigma_1^+\subset\Sigma_2^-\cup
  E_2$, et comme $\Sigma_1^+\cap E_2=\emptyset$, cela équivaut
  \struck{Si on avait} aussi à :
  \item[(2)] $\Sigma_1^+\subset\Sigma_2^-$
\end{itemize}
symétriquement, à :
\begin{itemize}
  \item[(2')] $\Sigma_2^+\subset\Sigma_1^-$
\end{itemize}
Comme $\Sigma_2^-\supset E_1$, \struck{vois.\ ouvert} i.e.\ $\Sigma_2^-$ est
un \emph{vois.}\ ouvert de $E_1$, \struck{il contient les} on a :
$\Sigma_2^-\cap\Sigma_1^+\neq\emptyset$ car
$\overline{\Sigma_1^+}\cap E_1\neq\emptyset$ (sinon $\Sigma_1^+$,
\add{ouvert,} serait fermé dans $X$, donc \struck{$\Sigma_1^+\neq\emptyset$}
$X$ serait disconnexe),
\marginal{$\overline{\Sigma_1^+}\cap E_1\neq\emptyset$ \ill{}}

\page{15}
et par suite $\Sigma_2^-\cup\Sigma_1^+$ serait un ouvert connexe de $X$ donc
de $X\setminus E_2$, contenant \add{la comp.\ conn.}\ $\Sigma_2^-$
\add{de $X\setminus E_2$}, donc \struck{\ill{}} égal à $\Sigma_2^-$, donc
$\Sigma_1^+\subset\Sigma_2^-$, cqfd

L'intersection des deux partitions
\[
X = \Sigma_1^+\sqcup E_1\sqcup\Sigma_1^-, \qquad
X = \Sigma_2^+\sqcup E_2\sqcup\Sigma_2^-
\]
de $X$ est
\[
X = \Sigma_1^+\sqcup E_1\sqcup\Sigma_{12}\sqcup E_2\sqcup\Sigma_2^+
\]
\note{une accolade au-dessus de $\Sigma_1^+\sqcup E_1\sqcup\Sigma_{12}$ porte
$\Sigma_2^-$, une accolade au-dessous de $\Sigma_{12}\sqcup E_2\sqcup
\Sigma_2^+$ porte $\Sigma_1^-$}
avec
\[
\Sigma_{12} \overset{\mathrm{déf}}{=} \Sigma_1^-\cap\Sigma_2^- .
\]

Dans cette partition, $\Sigma_1^+$, $\Sigma_{12}$, $\Sigma_2^+$ sont ouverts,
$E_1$, $E_2$ sont fermés, et on a
\[
\Sigma_1^+ \subsetneq \overline{\Sigma_1^+} \subset \Sigma_1^+\cup E_1
\]
\[
\overline{\Sigma_{12}} \subset \overline{\Sigma_1^-}\cap
\overline{\Sigma_2^-} \subset (\Sigma_1^-\cup E_1)\cap(\Sigma_2^-\cup E_2)
\subset \Sigma_{12}\cup E_1\cup E_2
\]
\[
\Sigma_2^+ \subsetneq \overline{\Sigma_2^+} \subset \Sigma_2^+\cup E_2
\]

Notons aussi que
\[
\overline{\Sigma_{12}}\cap E_1\neq\emptyset, \qquad
\overline{\Sigma_{12}}\cap E_2\neq\emptyset
\]
Car si on avait p.ex.\ $\overline{\Sigma_{12}}\cap E_1=\emptyset$, on
\struck{alors $\Sigma_{12}\sqcup E_2\sqcup\Sigma_2$} \add{en considérant}
$\overline{\Sigma_1^-}\cap E_1=\emptyset$, [car
\[
\overline{\Sigma_1^-} = \overline{\Sigma_2^+\sqcup E_2\sqcup\Sigma_{12}} =
\overline{\Sigma_2^+}\cup E_2\cup\overline{\Sigma_{12}} \subset
\Sigma_1^-\cup\overline{\Sigma_{12}}
\]
\note{sous $\overline{\Sigma_2^+}$, une accolade porte
$\subset\Sigma_2^+\cup E_2$}
ce qui n'est pas.

\page{16}
\emph{Cas particuliers}

Supposons que $\forall i\in I$, \struck{$E_i$ soit rare dans $X$} on ait (où
$\Sigma_i^+$ $\Sigma_i^-$ sont les deux comp.\ conn.\ de
$X\setminus\Sigma_i$) \struck{Alors on trouve que}
\note{$X\setminus\Sigma_i$ : sic, pour $X\setminus E_i$}
\[
\overline{\Sigma_i^+} = \Sigma_i^+\cup E_i, \qquad
\overline{\Sigma_i^-} = \Sigma_i^-\cup E_i .
\]
Alors je dis que
\[
\overline{\Sigma_{12}} = \Sigma_{12}\cup E_1\cup E_2
\]

\struck{Supposons} Soit $J\subset I$ tel que $I=J\sqcup\{i\}$. Considérons
la partition \add{$\Pi_J$} de $X$ \struck{définie} définie par les $E_j$
($j\in J$), \struck{intersection} \add{intersection} des partitions
\add{ternaires $\Pi_j$} de $X$ ($j\in J$) définies par les $E_j$ \ill{}
($j\in J$)
\[
X = E_j \sqcup \coprod_{\alpha\in\pi_0(X\setminus E_j)} \Sigma_j^{\alpha}
\]
— c'est aussi la partition définie par celle de $\bigcup_{j\in J}E_j$ en les
$E_j$, et la partition \add{$\Pi_J^0$} de $X\setminus\bigcup_{j\in J}E_j$
intersection des partitions binaires $\Pi_{J,j}^0$ \add{($j\in J$)}
(induites sur $U_J = X\setminus\bigcup_{j\in J}E_j$ par les $\Pi_j$
($j\in J$)). \uncertain{Pour}
\marginal{introduisons la partition
$\Pi_j = \{E_j, \Sigma_j^+, \Sigma_j^-\}$ \ill{}}

\page{17}
cette dernière, si $\operatorname{card}J = \operatorname{card}I - 1 = n-1$,
on verra (par récurrence) que le cardinal de l'ens.\ des classes est
$(n-1)+1=n$.

Ceci posé, on voit d'abord que $E_i$ est dans l'une \add{(unique)} des
classes, qui sont des ouverts. Soit $U_{J,i}$ la classe (ou “composante”)
qui contient $E_i$, et considérons les deux composantes $\Sigma_i^+$,
$\Sigma_i^-$ de $X\setminus E_i$.

\emph{Lemme} Toute classe de $\Pi_J^0$, distincte de $U_{J,i}$, est
contenue dans $\Sigma_i^+$ ou $\Sigma_i^-$.

C'est presque une tautologie : Dire \add{qu'une} \struck{la} classe donnée
\add{$U_\alpha$ de $\Pi_J^0$} \struck{un} ne contient pas $E_i$ i.e.\ est
$\neq U_{J,i}$, signifie qu'il existe un $j\in J$ tel que la
$\Sigma_j^+\in\pi_0(X\setminus E_j)$ \struck{choisie} “composante” de cette
classe \struck{est} soit l'élément de $\pi_0$ qui ne contient pas $E_i$ ; il
est donc contenu dans l'un des éléments \add{(soit $\Sigma_i^-$)} de
$\pi_0(X\setminus\Sigma_i)$, d'après l'étude du cas
$\operatorname{card}I=2$, \struck{et} à fortiori
$U_\alpha\subset\Sigma_i^-$, OK.
\note{$X\setminus\Sigma_i$ : sic, pour $X\setminus E_i$}

\page{18}
Dém.\ par récurrence, c'est trivial si $\operatorname{card}I=0,1$, le cas
$\operatorname{card}I=2$ a été vu. Cas général
$\operatorname{card}I\geq 3$, on choisit $J$ comme dessus. Il suffit de voir
que les ens.
\[
U_{J,i}\cap\Sigma_i^+, \qquad U_{J,i}\cap\Sigma_i^-
\]
ont la forme précédente, il suffit de le voir pour l'un des deux. Par hyp.\ de récurrence, on a
\[
U_{J,i} = X\setminus\coprod_{\alpha\in K_0}\widehat{\Sigma}_\alpha^+
\]
où
\[
\begin{cases}
K_0\subset J\\
\Sigma_\alpha^+\in\pi_0(X\setminus E_\alpha) \quad (\forall\alpha\in K_0)\\
\widehat{\Sigma}_\alpha^+ = \Sigma_\alpha^+\cup E_\alpha .
\end{cases}
\]
\note{« la forme précédente » est celle que donne le lemme de la page 19}

\struck{On} \uncertain{On a} donc
\[
U = U_{J,i}\cap\Sigma_i^- = X\setminus\Bigl(\coprod_{\alpha\in K_0}
\widehat{\Sigma}_\alpha^+ \cup \widehat{\Sigma}_i^+\Bigr)
\]
\struck{Soit} \uncertain{Au} \uncertain{dernier} $=$, dans la
\struck{\ill{}} réunion $\bigcup_{\alpha\in K_0}\widehat{\Sigma}_\alpha^+$,
il suffit de se borner : \ill{} sur les \struck{$\Sigma$}
$\alpha\in K_0$ tels que $\widehat{\Sigma}_\alpha^+\not\subset
\widehat{\Sigma}_i^+$ i.e.\ $\Sigma_\alpha^+\not\subset\Sigma_i^+$ soit
$K_1$ cet ens., donc
\[
U = X\setminus\Bigl(\bigcup_{\alpha\in K_1}\widehat{\Sigma}_\alpha^+
\cup\widehat{\Sigma}_i^+\Bigr)
\]
\note{la démonstration reprend en haut de la page 20}

\page{19}
Ainsi les classes de $\Pi_I^0$ sont les classes de $\Pi_J^0$ distinctes de
$U_{J,i}$, et de plus les deux parmi
\[
U_{J,i}\cap\Sigma_i^+, \qquad U_{J,i}\cap\Sigma_i^-
\]
qui sont non vides. Or les deux classes en question sont non vides, car
\struck{\ill{}} $U_{J,i}$ est un voisinage de $E_i$ et
$\overline{\Sigma_i^+}\cap E_i\neq\emptyset$,
$\overline{\Sigma_i^-}\cap E_i\neq\emptyset$. On voit donc que le nb des
classes est \struck{\ill{}}
$(n-1)+2 = n+1 = \operatorname{card}I+1$ (ce qui achève aussi la dém.\ par
récurrence sur le nb des classes).
\note{ce paragraphe prolonge le raisonnement de la page 17 ; le lemme qui
suit est celui dont les pages 18 et 20 donnent la démonstration}

\emph{Lemme} Toute classe de $\Pi_I^0$ est de la forme
\[
X\setminus\bigcup_{\alpha\in K}(\Sigma_\alpha^+\cup E_\alpha)
\]
\note{au-dessus de $\Sigma_\alpha^+\cup E_\alpha$, il écrit
$\widehat{\Sigma}_\alpha^+$}
où $K$ est une partie \add{\uncertain{non vide}} de $I$, et
$\forall\alpha\in K$, \struck{\ill{}}
$\Sigma_\alpha^+\in\pi_0(X\setminus E_\alpha)$, tels que
$\forall\alpha,\beta\in K$, $\alpha\neq\beta$, on ait
$\widehat{\Sigma}_\alpha^+\cap\widehat{\Sigma}_\beta^+=\emptyset$
($\Longleftrightarrow \Sigma_\alpha^+\cap\Sigma_\beta^+=\emptyset$).
\marginal{$K$ est $\neq\emptyset$ si $I\neq\emptyset$}

\page{20}
Soit $\alpha_0\in K_1$, je dis que \uncertain{on a} :
\[
\widehat{\Sigma}_{\alpha_0}\cap\widehat{\Sigma}_i^+ = \emptyset \quad
\text{i.e.}\quad \Sigma_{\alpha_0}\cap\Sigma_i = \emptyset
\]
(ce qui achèvera la démonstration, car on aura
\[
U = X\setminus\bigcup_{\alpha\in K}\widehat{\Sigma}_\alpha^+ \qquad
\bigl|\ K = K_1\cup\{i\}\ \struck{\ill{}}
\]
comme dans le lemme.

Or \struck{\ill{}} \ill{} :
\[
E_i\subset U_{J,i} = X\setminus\bigcup_{\alpha\in K_0}
\widehat{\Sigma}_\alpha^+
\]
ce qui implique que
\[
E_i\subset X\setminus\widehat{\Sigma}_{\alpha_0}^+ = \Sigma_{\alpha_0}^-
\]
i.e.
\note{à droite, un schéma vertical : de haut en bas $\Sigma_i^+$ ou
$\Sigma_i^-$, $E_i$, $\Sigma_\alpha$, $E_\alpha$, $\Sigma_{\alpha_0}^+$,
réunis par une accolade}
d'autre part $\alpha_0\in K_1$ signifie
\[
\Sigma_{\alpha_0}^+\not\subset\Sigma_i^+
\]
ce qui \struck{implique} signifie aussi que
$\Sigma_i^+\cap\Sigma_\alpha^+=\emptyset$, cqfd.

\emph{Corollaire du lemme} Pour une classe $V$ de $\Pi_I^0$, son
expression
\[
V = X\setminus\bigcup_{\alpha\in K}\widehat{\Sigma}_\alpha^+
\]
du lemme précédent est \emph{unique} — on a
\[
\begin{cases}
K = \{\alpha\in I \mid \overline{V}\cap E_\alpha\neq\emptyset \ \text{ i.e. }\ E_\alpha\cap\dot{V}\neq\emptyset\}\\
\Sigma_\alpha^+ \text{ est l'unique comp.\ de } X\setminus E_\alpha
\text{ contenue dans } X\setminus V
\end{cases}
\]
\note{le corollaire, encadré d'un trait vertical, s'arrête au bas du
feuillet ; la suite est au lot suivant}

\end{document}
